\documentclass[10pt,aspectratio=169,serif,fleqn]{beamer}

\usepackage[fontset=none, scheme=plain]{ctex}
\IfFontExistsTF{Songti SC}{%
	\setCJKmainfont{Songti SC}[AutoFakeSlant=0.2]%
	\setCJKsansfont{Songti SC}[AutoFakeSlant=0.2]%
}{%
	\setCJKmainfont{FandolSong-Regular}[Extension=.otf,BoldFont=FandolSong-Bold,AutoFakeSlant=0.2]%
	\setCJKsansfont{FandolSong-Regular}[Extension=.otf,BoldFont=FandolSong-Bold]%
}
\usepackage{amsmath,amssymb,graphicx,tikz,booktabs,multicol,calligra}
\usepackage{pgfplots}
\pgfplotsset{compat=1.18}
\usepackage{tcolorbox}
\tcbuselibrary{skins}
\usepackage{XyygU}
\usepackage{caption}
\captionsetup[figure]{name=图,labelsep=period}

\usetikzlibrary{arrows.meta,positioning,calc}

\setlength{\abovedisplayskip}{0pt plus 2pt minus 2pt}
\setlength{\belowdisplayskip}{0pt plus 2pt minus 2pt}

\title[函数的应用(一)练习]{3.4 函数的应用（一）（练习）}
\subtitle{函数模型的实际应用}
\author{}
\institute{}
\date{2026-10-30}

\begin{document}

\begin{frame}
  \titlepage
\end{frame}

% ========== 第1题 ==========
\begin{frame}[t]{第1题：汽车刹车距离}
\textbf{1.} 若用模型 $y=ax^2$ 描述汽车紧急刹车后滑行的距离 $y$（单位：m）与刹车时的速率 $x$（单位：km/h）的关系，而某种型号的汽车在速率为 $60$ km/h 时，紧急刹车后滑行的距离为 $20$ m。在限速为 $100$ km/h 的高速公路上，一辆这种型号的车紧急刹车后滑行的距离为 $50$ m，那么这辆车是否超速行驶？
\end{frame}

\begin{frame}[t]{第1题 解答}
	\vspace{-2em}
\begin{tcolorbox}[colback=orange!10,colframe=orange,title=\textbf{解}]
将 $x=60$，$y=20$ 代入 $y=ax^2$：
\[
20=a\cdot 60^2=3600a,\quad a=\frac{20}{3600}=\frac1{180}.
\]
故模型为 $y=\dfrac{x^2}{180}$。

\vspace{0.5em}
当 $y=50$ 时：
\[
50=\frac{x^2}{180},\quad x^2=9000,\quad x=\sqrt{9000}\approx 94.9\ \text{km/h}.
\]

\vspace{0.5em}
因为 $94.9<100$，所以这辆车 \alert{没有超速}行驶。
\end{tcolorbox}
\end{frame}

% ========== 第2题 ==========
\begin{frame}[t]{第2题： 某广告公司要为客户设计一幅周长为 $l$（单位：m）的矩形广告牌，如何设计这个广告牌可以使广告牌的面积最大？}


\vspace{-1em}
\begin{tcolorbox}[colback=orange!10,colframe=orange,title=\textbf{解}]
	设矩形一边长为 $x$，则另一边长为 $\dfrac{l}{2}-x$，面积：
	\[
	S=x\left(\frac{l}{2}-x\right)=-x^2+\frac{l}{2}x=-\left(x-\frac{l}{4}\right)^2+\frac{l^2}{16}.
	\]
	当 $x=\dfrac{l}{4}$ 时，$S$ 取最大值 $\dfrac{l^2}{16}$。
	
	\vspace{0.5em}
	此时矩形两边均为 $\dfrac{l}{4}$，即广告牌设计为 \alert{边长为 $\dfrac{l}{4}$ 的正方形}时面积最大。
\end{tcolorbox}
\end{frame}



% ========== 第3题 ==========
\begin{frame}[t]{第3题：产品成本与利润}
\textbf{3.} 某公司生产某种产品的固定成本为 $150$ 万元，而每件产品的可变成本为 $2500$ 元，每件产品的售价为 $3500$ 元。若该公司所生产的产品全部销售出去，则

(1) 设总成本为 $y_1$（单位：万元），单位成本为 $y_2$（单位：万元），销售总收入为 $y_3$（单位：万元），总利润为 $y_4$（单位：万元），分别求出它们关于总产量 $x$（单位：件）的函数解析式；

(2) 根据所求函数的图象，对这个公司的经济效益做出简单分析。
\end{frame}

\begin{frame}[t]
	\vspace{-2em}
\begin{tcolorbox}[colback=orange!10,colframe=orange,title=\textbf{解：(1)}]
$2500$ 元 $=0.25$ 万元，$3500$ 元 $=0.35$ 万元。

\vspace{0.3em}
总成本：
\[
y_1=150+0.25x\quad (\text{万元});
\]
单位成本：
\[
y_2=\frac{y_1}{x}=\frac{150}{x}+0.25\quad (\text{万元/件});
\]
销售收入：
\[
y_3=0.35x\quad (\text{万元});
\]
总利润：
\[
y_4=y_3-y_1=0.35x-(150+0.25x)=0.1x-150\quad (\text{万元}).
\]
\end{tcolorbox}
\end{frame}

\begin{frame}[t]{第3题 解答：(2)}
\begin{tcolorbox}[colback=orange!10,colframe=orange,title=\textbf{解：(2)}]
令 $y_4=0$，得 $0.1x-150=0$，$x=1500$（件）。

\vspace{0.5em}
\textbf{盈亏分析：}
\begin{itemize}
  \item 当 $x<1500$ 件时，$y_4<0$，公司亏损；
  \item 当 $x=1500$ 件时，$y_4=0$，盈亏平衡（保本点）；
  \item 当 $x>1500$ 件时，$y_4>0$，公司盈利，且产量越大利润越高。
\end{itemize}

\vspace{0.5em}
因此该公司应争取产量超过 $1500$ 件才能获利。
\end{tcolorbox}
\end{frame}

\begin{frame}[c]
  \begin{center}
    {\Huge\calligra Thanks!}
  \end{center}
\end{frame}

\end{document}
